% TOP_PROBLEM: {"id": 3157, "title": "Subgraph of large average degree and large girth.", "category_id": 3, "category": {"id": 3, "name": "graph_theory", "display_name": "Graph Theory", "description": "Problems involving graphs, networks, and their properties.", "slug": "graph-theory", "order_index": 3, "created_at": "2026-07-31T15:26:25.671Z"}, "status": "open", "problem_number": "OPG-46708", "set_id": 12}
% FIRST_POSTED: 2026-10-01
% ATTEMPT_STATUS: unresolved
% UnsolvedMath ID 3157; source catalogue code OPG-46708
% !TeX program = lualatex
% GPT-6 Astra Ultra research attempt; independent partial work, 2026-10-01.
\documentclass[11pt,a4paper]{article}
\usepackage[margin=25mm]{geometry}
\usepackage{fontspec}
\setmainfont{lmroman10-regular.otf}[BoldFont=lmroman10-bold.otf,
 ItalicFont=lmroman10-italic.otf,BoldItalicFont=lmroman10-bolditalic.otf]
\usepackage{amsmath,amssymb,mathtools}
\usepackage{unicode-math}
\setmathfont{latinmodern-math.otf}
\AtBeginDocument{\let\setminus\smallsetminus}
\usepackage{xurl,hyperref}
\hypersetup{colorlinks=true,urlcolor=blue}
\setcounter{secnumdepth}{0}
\setlength{\parindent}{0pt}
\setlength{\parskip}{5pt}
\setlength{\emergencystretch}{3em}
\begin{document}
\section*{Subgraph of large average degree and large girth.}\label{3157}
{\small UnsolvedMath ID 3157 · OPG-46708 · Graph Theory · OpenGarden}
\subsection{Definitions and mathematical statement}
For a nonempty finite simple graph $G=(V,E)$, its average degree
is $\overline d(G)=2|E|/|V|$. Its girth is the minimum length of
a simple cycle; set the girth to $+\infty$ when the graph is a
forest. A subgraph may be obtained by deleting vertices and
edges, so it need not be induced or spanning.

The conjecture has the following quantifier order:
\[
 \forall g,k\in\mathbb Z_{\geq1}\quad
 \exists d=d(g,k)\in\mathbb Z_{\geq1}\quad
 \forall G\quad
 \overline d(G)\geq d\Longrightarrow
 \exists H\subseteq G:
 \overline d(H)\geq k,\quad\operatorname{girth}(H)>g.
\]
Here $G$ and the selected $H$ are nonempty finite simple graphs.
The threshold is independent of the order and maximum degree
of $G$. The conclusion means that $H$ contains no cycle with at
most $g$ edges while retaining the prescribed average degree.

No regularity assumption is imposed on $G$. In particular, its
large average degree may coexist with widely varying vertex
degrees. One may discard edges of short cycles and choose a
smaller vertex set, but both resulting conditions must hold in
the same subgraph. The statement demands an existence threshold
for each pair $(g,k)$; it does not prescribe a numerical formula
or a polynomial bound for $d(g,k)$. A counterexample would fix
some pair $(g,k)$ and produce graphs of unbounded average degree
with no suitable subgraph.
\subsection{Short English statement}
Does sufficiently large average degree force a subgraph with any prescribed average degree and no short cycles?
\subsection{Research attempt: alteration with a controlled maximum degree}
\paragraph{Scope, context, and attack.}
GPT-6 Astra Ultra, 1 October 2026. The threshold must depend only on
$g,k$, not on the size or maximum degree of the input. The primary
paper [S3] identifies this as Thomassen's conjecture and develops the
$C_4$-free case. The elementary sampling argument already outlined in [S3] is recovered
below with explicit constants. It proves an arbitrary-girth statement
under the additional hypothesis that maximum degree is at most
$K$ times average degree, for a fixed $K$. It is not a new resolution
of the unrestricted conjecture. The strategy is independent edge sampling
followed by deleting one edge from each surviving short cycle.

\paragraph{1. Small girth requirements.}
For $g\le2$, simplicity already excludes cycles of length at most $g$,
so $H=G$ works when the average degree is at least $k$. For $g=3$,
a random bipartition places each edge across the partition with
probability $1/2$. Some bipartite spanning subgraph therefore retains
at least half the edges, has average degree at least
$\overline d(G)/2$, and has no triangle. Thus $d(3,k)=2k$ is sufficient.
This argument controls triangles without any maximum-degree assumption.
The difficulty begins when one must eliminate even short cycles as well.

\paragraph{2. An explicit bound for graphs with bounded degree ratio.}
Let $n=|V(G)|$, $d=\overline d(G)$, and $\Delta=\Delta(G)$, and
assume $d\ge1$ and $\Delta\le Kd$, where $K\ge1$ is fixed. Fix
$g\ge3$ and retain each edge independently with probability
\[
 p=\tfrac12\Delta^{-1+1/g}.
\]
There are at most $n\Delta^{\ell-1}/(2\ell)$ simple $\ell$-cycles:
choose a starting vertex and the next $\ell-1$ vertices of an oriented
walk, then require its last vertex to be adjacent to the first;
each actual cycle is counted $2\ell$ times. If $Z$ counts surviving
cycles of lengths $3,\ldots,g$, linearity of expectation gives
\[
 \mathbb E Z\le n\sum_{\ell=3}^g
 \frac{2^{-\ell}}{2\ell}\Delta^{\ell/g-1}
 \le \frac n6\sum_{\ell=3}^{\infty}2^{-\ell}=\frac n{24}.
\]
The expected number $M$ of sampled edges satisfies
\[
 \mathbb E M=\frac{npd}{2}
 \ge\frac{n\Delta^{1/g}}{4K}.
\]
Choose an outcome with $M-Z\ge\mathbb E(M-Z)$. List its short
cycles and remove one edge from each cycle still present when it is
processed. At most $Z$ edges are removed. Edge deletion creates no
new cycle, so the resulting spanning subgraph $H$ has girth greater
than $g$ and
\[
 \overline d(H)\ge\frac{\Delta^{1/g}}{2K}-\frac1{12}.
\]
In particular, $d\ge(2K(k+1))^g$ suffices, since $\Delta\ge d$.
Isolated vertices may be retained in this calculation; deleting them
can only increase average degree. For regular graphs this gives the
explicit sufficient threshold $(2(k+1))^g$.

\paragraph{Verification and obstruction to removing the hypothesis.}
The cycle count is an upper bound even when most walks are not cycles;
independence is used only for the probability $p^\ell$ that all edges
of one fixed cycle survive. No independence between different cycle
events is asserted. Choosing an outcome by the expectation of $M-Z$
also avoids incorrectly optimizing the two quantities separately.

The bound contains $K$ essentially in this proof. In $K_{t,t^2}$,
$\Delta=t^2$ and $d=2t^2/(t+1)$, so $\Delta/d=(t+1)/2$ is unbounded.
These graphs are not counterexamples to the conjecture; they show that
high average degree alone does not meet the extra input condition.
A reduction extracting a sufficiently dense subgraph with a degree
ratio bounded solely in terms of the desired parameters would close
this particular gap, but no such reduction is established here.

\subsection{Final status}
UNRESOLVED. The triangle-free case and a quantitative bounded-degree-ratio
case are proved. The unrestricted irregular case remains beyond this
attempt. Completion estimate: no defensible global percentage.

\subsection{Sources}
\begin{itemize}
\item[S1] ULAM.AI and the UnsolvedMath Contributors, supplied problems.json, record 3157 (OPG-46708), OpenGarden collection. Catalogue identity and source transcription; CC BY 4.0.
\url{https://huggingface.co/datasets/ulamai/UnsolvedMath}.
\item[S2] Open Problem Garden, Subgraph of large average degree and large girth.. Original problem and discussion; the mathematical formulation below is expanded from the supplied source record.
\url{http://www.openproblemgarden.org/op/subgraph_of_large_average_degree_and_large_average_degree}.
\item[S3] R. Montgomery, A. Pokrovskiy and B. Sudakov, \emph{$C_4$-free subgraphs with large average degree}, Israel Journal of Mathematics 246 (2021), 55--71; author manuscript checked 1 October 2026.
\url{https://people.math.ethz.ch/~sudakovb/C_4-free-subgraphs.pdf}.
\end{itemize}
\end{document}
